\chapter{Universal Properties}

\begin{overviewbox}
Category theory characterizes objects by their relationships to other
objects, not by what they are made of inside. The cleanest expression of this
principle is the \emph{universal property}: a description of an object that
says, in effect, ``this object is the unique one (up to unique isomorphism)
that participates in such-and-such a way in such-and-such a situation.''

This chapter introduces the universal-property pattern through its two
simplest cases: \emph{initial objects} (those with a unique map out to every
other object) and \emph{terminal objects} (those with a unique map in from
every other object). We then state the general pattern, prove that universal
objects are unique up to unique isomorphism, and use the pattern to preview
\emph{products} and \emph{coproducts} --- the first universal constructions
involving more than one object at a time.

Universal properties are everywhere from this chapter onward. The investment
made here in understanding the pattern pays off repeatedly.
\end{overviewbox}


% ═══════════════════════════════════════════════════════════════════════════
\begin{nameorigin}
\term{Universal property} captures one of the deepest patterns in
mathematics: characterizing an object not by what it \emph{is} but by
what it \emph{does} — its universal mapping behavior. The terminology
is borrowed from physics and analysis (universal cover, universal
constant, etc.) where ``universal'' means ``applies in every case.''
In category theory, a universal property is an object that
``solves a problem uniquely up to unique isomorphism'' — it is the
universal solution to a mapping question. The pattern was formalized
by Bourbaki (mid-20th century) and made into the central technique of
category theory by Mac Lane and Eilenberg.

\term{Initial} and \term{terminal} objects are the universal-property
basics: initial = unique map \emph{out} to everything; terminal =
unique map \emph{in} from everything. ``Initial'' and ``terminal'' are
from the structure of arrows — first/last in any path.
\end{nameorigin}

\section{Initial Objects}
\label{sec:initial}
% ═══════════════════════════════════════════════════════════════════════════

\subsection{Informally: One Way Out}

An \term{initial object} in a category is an object from which there is
exactly one morphism to every other object --- including itself. ``Exactly
one'' is the key phrase: not zero, not two, not many. The initial object
has a unique outbound connection to every destination.

If you imagine the category as a transit map, an initial object is a station
from which there is one and only one route to every station in the city.
There may be many such routes from other stations, but starting at the
initial object, your route is forced.

\subsection{Expanding: Looking for Initial Objects}

\begin{example}{$\mathbf{0}$ Is Initial in Itself}
The empty category $\mathbf{0}$ has no objects. The statement ``for every
object $A$, there is exactly one morphism from $X$ to $A$'' is vacuously
true for any $X$, but $\mathbf{0}$ has no objects to play the role of $X$
either. Vacuously, the empty category is its own (trivial) example.
\end{example}

\begin{example}{$\mathbf{1}$ Is Initial in $\mathbf{2}$ Read Backwards?}
Recall $\mathbf{2}$: two objects $0, 1$ with one non-identity morphism
$\iota : 0 \to 1$. The object $0$ has:
\begin{align*}
  \mathrm{Hom}(0, 0) &= \{\mathrm{id}_0\}, \\
  \mathrm{Hom}(0, 1) &= \{\iota\}.
\end{align*}
Exactly one morphism to each object. So $0$ is the initial object of
$\mathbf{2}$. The object $1$, by contrast, has $\mathrm{Hom}(1, 0) =
\varnothing$ --- no morphism to $0$ at all. So $1$ is not initial.
\end{example}

\begin{example}{The Recipe Category}
The recipe category has Chop, Sauté, Add with morphisms
$\mathrm{Chop} \to \mathrm{Saut\acute{e}} \to \mathrm{Add}$ (and their
composite). From Chop:
\begin{align*}
  \mathrm{Hom}(\mathrm{Chop}, \mathrm{Chop}) &= \{\mathrm{id}_{\mathrm{Chop}}\}, \\
  \mathrm{Hom}(\mathrm{Chop}, \mathrm{Saut\acute{e}}) &= \{f\}, \\
  \mathrm{Hom}(\mathrm{Chop}, \mathrm{Add}) &= \{g \circ f\}.
\end{align*}
Exactly one morphism to each. So Chop is initial. It is the unique step
from which a unique route reaches every other step.
\end{example}

\begin{example}{No Initial Object}
Take the indiscrete category on two objects $\{A, B\}$: hom-sets are
all singletons including $\mathrm{Hom}(A, B)$, $\mathrm{Hom}(B, A)$,
$\mathrm{Hom}(A, A) = \{\mathrm{id}_A\}$, $\mathrm{Hom}(B, B) =
\{\mathrm{id}_B\}$. Both $A$ and $B$ have exactly one morphism to each
object, so both are initial. We will see in the next section that this
forces $A \cong B$, but here we note: \emph{there can be more than one
initial object} in a category, but they must be isomorphic.
\end{example}

\begin{dialog}
To check whether an object $X$ is initial, I ask, for each object $A$ in the
category: how many morphisms $X \to A$ are there?
\begin{itemize}
  \item If exactly one for every $A$: $X$ is initial.
  \item If zero for some $A$: $X$ is not initial.
  \item If two or more for some $A$: $X$ is not initial.
\end{itemize}
The check is local at $X$ but quantifies over every $A$ in the category. In
small categories, this is a finite check. In large categories, it is usually
a structural argument.
\end{dialog}

\subsection{Formalizing: The Definition}

An object $X \in \mathrm{Ob}(\mathcal{C})$ is \term{initial} if for every
object $A \in \mathrm{Ob}(\mathcal{C})$, the hom-set $\mathrm{Hom}(X, A)$
is a singleton:
\begin{equation*}
  | \mathrm{Hom}_{\mathcal{C}}(X, A) | = 1
  \quad \text{for every } A \in \mathrm{Ob}(\mathcal{C}).
\end{equation*}
We write $\mathbf{0}$ for an initial object (when one exists), and $!_A$
or simply $!$ for the unique morphism $\mathbf{0} \to A$.

\begin{howtosay}
\begin{itemize}
  \item ``initial object'' is sometimes called the \emph{universal repelling
        object}, or the \emph{strict initial} in some refined contexts.
  \item The notation $\mathbf{0}$ (boldface zero) is conventional. Some
        authors use $\bot$ (read ``bottom''), particularly when the category
        comes from logic.
  \item The unique morphism is often written $!_A$ (``shriek to $A$'') or
        $!_{\mathbf{0}, A}$. The exclamation point itself is read
        ``shriek'' or ``bang.''
  \item Drawn as a dashed arrow when emphasizing uniqueness:
        $\mathbf{0} \dashrightarrow A$.
\end{itemize}
\end{howtosay}

\subsection{Reorganizing: The Examples Restated}

\begin{align*}
  &\text{$\mathbf{0}$ in $\mathbf{2}$:}
    && |\mathrm{Hom}(0, 0)| = 1,\; |\mathrm{Hom}(0, 1)| = 1. \;\text{Initial.} \\
  &\text{Chop in Recipe:}
    && |\mathrm{Hom}(\mathrm{Chop}, X)| = 1 \;\text{for } X \in \{\mathrm{Chop}, \mathrm{Saut\acute{e}}, \mathrm{Add}\}. \\
  &\text{Indiscrete on $\{A, B\}$:}
    && \text{both $A$ and $B$ initial; both isomorphic.} \\
  &\text{$\mathbf{0}$ (empty category):}
    && \text{vacuously its own (trivial) initial object.}
\end{align*}

\subsection{Simplifying: Initial in One Line}

\begin{equation*}
  X \text{ is initial in } \mathcal{C}
  \;\iff\;
  \forall\, A \in \mathcal{C},\; \exists !\, f \in \mathrm{Hom}(X, A).
\end{equation*}

The symbol $\exists !$ is read ``there exists a unique.''

\subsection{Formally: Uniqueness Sketch}

\formalnote{Initial objects are unique up to unique isomorphism}
Suppose $X$ and $Y$ are both initial in $\mathcal{C}$. There is a unique
$f : X \to Y$ (because $X$ is initial), and a unique $g : Y \to X$ (because
$Y$ is initial). Consider $g \circ f : X \to X$. There is exactly one
morphism $X \to X$ (because $X$ is initial), and $\mathrm{id}_X$ is such a
morphism. So $g \circ f = \mathrm{id}_X$. By the same argument,
$f \circ g = \mathrm{id}_Y$. Hence $f$ is an isomorphism, and uniqueness of
$f$ as a morphism gives uniqueness of the isomorphism.

We will state and reuse this argument in its general form in
\S5.3. Notice that no inspection of what $X$ or $Y$ are ``made of'' is
involved --- the entire proof uses only the universal property.


% ═══════════════════════════════════════════════════════════════════════════
\section{Terminal Objects}
\label{sec:terminal}
% ═══════════════════════════════════════════════════════════════════════════

\subsection{Informally: One Way In}

A \term{terminal object} is the mirror image of an initial object: an object
to which there is exactly one morphism from every other object. Instead of
``one way out,'' it is ``one way in.''

\subsection{Expanding: A Categorical Mirror}

\begin{example}{$1$ Is Terminal in $\mathbf{2}$}
In $\mathbf{2}$, the object $1$ has:
\begin{align*}
  \mathrm{Hom}(0, 1) &= \{\iota\}, \\
  \mathrm{Hom}(1, 1) &= \{\mathrm{id}_1\}.
\end{align*}
Exactly one morphism into $1$ from every object. So $1$ is terminal. (The
object $0$ is not terminal, since $\mathrm{Hom}(1, 0) = \varnothing$.)
\end{example}

\begin{example}{Add in the Recipe Category}
In the recipe category, Add receives the morphisms $\mathrm{id}_{\mathrm{Add}}$,
$g$, and $g \circ f$. One from each step. So Add is terminal.
\end{example}

\begin{example}{An Object That Is Both}
In $\mathbf{1}$ (one object, one morphism), the unique object $*$ has
exactly one morphism out (to itself) and exactly one morphism in (from
itself). It is simultaneously initial and terminal. Such an object is called
a \term{zero object}.
\end{example}

\subsection{Formalizing: Duality}

An object $X \in \mathrm{Ob}(\mathcal{C})$ is \term{terminal} if for every
object $A \in \mathrm{Ob}(\mathcal{C})$, the hom-set $\mathrm{Hom}(A, X)$
is a singleton:
\begin{equation*}
  | \mathrm{Hom}_{\mathcal{C}}(A, X) | = 1
  \quad \text{for every } A \in \mathrm{Ob}(\mathcal{C}).
\end{equation*}

\textbf{Duality.} An object is terminal in $\mathcal{C}$ iff it is initial
in $\mathcal{C}^{\mathrm{op}}$. The two notions are dual to each other in
the formal sense of \S2.4: any statement about initial objects becomes a
statement about terminal objects when you replace $\mathcal{C}$ by
$\mathcal{C}^{\mathrm{op}}$, and vice versa.

\begin{howtosay}
\begin{itemize}
  \item ``terminal object'' is also called the \emph{universal attracting
        object}, or in some texts the \emph{final object}. ``Terminal'' is
        most common in modern usage.
  \item Notation: $\mathbf{1}$ (boldface one) is standard. Some authors
        use $\top$ (read ``top''), especially in logic contexts.
  \item The unique morphism into a terminal object is also written $!$ or
        $!_A : A \to \mathbf{1}$.
  \item A \emph{zero object}, which is both initial and terminal, is
        sometimes denoted $0$, $\mathbf{0}$, or $*$.
\end{itemize}
\end{howtosay}

\subsection{Reorganizing: The Same Theorem, Twice}

The uniqueness-up-to-unique-isomorphism result we proved for initial objects
holds, by duality, for terminal objects: any two terminal objects are
isomorphic via a unique isomorphism. This is not a second proof to write;
it is the same proof, transported across the opposite-category construction.

This is the first concrete payoff of \S2.4. We proved one result and got
two theorems. From here on, every result about a universal property comes
in two flavors (covariant and contravariant), and we will prove only one
of them.

\subsection{Simplifying: Initial and Terminal}

\begin{align*}
  X \text{ initial}
    &\iff \forall A,\; \exists ! \, f : X \to A. \\
  X \text{ terminal}
    &\iff \forall A,\; \exists ! \, f : A \to X.
\end{align*}

\subsection{Formally: Zero Morphisms}

\formalnote{Zero objects and zero morphisms}
If $\mathcal{C}$ has a zero object $\mathbf{0}$ (an object that is both
initial and terminal), then for any pair $A, B$ there is a distinguished
morphism
\begin{equation*}
  0_{A,B} : A \to B,
\end{equation*}
obtained by composing the unique $A \to \mathbf{0}$ with the unique
$\mathbf{0} \to B$. This is the \term{zero morphism} from $A$ to $B$.
Zero morphisms have the property that any composite involving a zero
morphism is again a zero morphism. They behave like the algebraic zero in
contexts where one applies.

\formalnote{Existence is not guaranteed}
Many categories have no initial or terminal object at all. The empty
category $\mathbf{0}$ has no terminal object (no object exists to receive
the unique morphism from elsewhere). The category of nonempty preorders
under monotone maps has neither. Existence of universal objects is
context-dependent.


% ═══════════════════════════════════════════════════════════════════════════
\section{The Universal Property Pattern}
\label{sec:universal-pattern}
% ═══════════════════════════════════════════════════════════════════════════

\subsection{Informally: ``Unique So That\ldots''}

The defining feature of every universal property is a phrase of the form:
\begin{quote}
  \emph{for every} object $X$ with property $P$, \emph{there exists a unique}
  morphism $u : X \to U$ (or $U \to X$) \emph{such that} some diagram
  commutes.
\end{quote}
The object $U$ is the universal one. The phrase ``unique so that'' is the
signature.

Initial and terminal objects fit this template in the simplest possible way:
$P$ is trivial (``$X$ is any object''), and the diagram has only one path.
More involved universal properties add data and conditions, but the shape
``for every \ldots, there exists a unique \ldots such that \ldots'' is
preserved.

\subsection{Expanding: A Worked-Out Template}

\begin{example}{Initial Object as a Universal Property}
\begin{quote}
  An object $\mathbf{0}$ is \emph{initial} if:\\
  \emph{for every} object $A$,\\
  \emph{there exists a unique} morphism $!_A : \mathbf{0} \to A$.
\end{quote}

No diagram is needed beyond the morphism itself, because there is no other
data to compare against.
\end{example}

\begin{example}{Terminal Object as a Universal Property}
\begin{quote}
  An object $\mathbf{1}$ is \emph{terminal} if:\\
  \emph{for every} object $A$,\\
  \emph{there exists a unique} morphism $!_A : A \to \mathbf{1}$.
\end{quote}
\end{example}

\begin{example}{The Universal Property of a Product (Preview)}
For objects $A$ and $B$ in $\mathcal{C}$, a \emph{product} is an object $P$
equipped with two morphisms $\pi_1 : P \to A$ and $\pi_2 : P \to B$ such that:
\begin{quote}
  \emph{for every} object $X$ and every pair of morphisms $f : X \to A$,
  $g : X \to B$,\\
  \emph{there exists a unique} morphism $\langle f, g \rangle : X \to P$
  \emph{such that} the diagram
  \medskip
  \begin{center}
  \begin{tikzcd}[column sep=3em, row sep=2em]
    & X
      \arrow[dl, "f"', bend right=15]
      \arrow[dr, "g", bend left=15]
      \arrow[d, dashed, "{\langle f, g \rangle}"] & \\
    A & P \arrow[l, "\pi_1"] \arrow[r, "\pi_2"'] & B
  \end{tikzcd}
  \end{center}
  commutes.
\end{quote}

Here the universal property has additional data ($\pi_1, \pi_2$) and a
diagram with multiple paths. But the signature phrase is the same.
\end{example}

\subsection{Formalizing: The Pattern}

A \term{universal property} for a class of objects is a specification of:
\begin{enumerate}
  \item A category $\mathcal{C}$ in which we look for the universal object.
  \item Some \term{data} $D$ that the universal object must carry.
  \item A class of \term{competitors}: pairs $(X, d_X)$ where $X$ is an
        object and $d_X$ is the same kind of data attached to $X$.
  \item A \term{uniqueness requirement}: for every competitor $(X, d_X)$,
        there is exactly one morphism from $X$ to the universal object (or
        from the universal object to $X$) compatible with the data.
\end{enumerate}

The universal object is the ``optimal'' member of the class of competitors
in a precise sense made by the uniqueness requirement.

\subsection{Reorganizing: Recognizing the Pattern}

Whenever you see a definition of the form
\begin{quote}
  \textit{For every $X$ with $\ldots$, there exists a unique $u : X \to U$
  such that $\ldots$ commutes,}
\end{quote}
you are looking at a universal property. The next few chapters introduce
many such definitions:
\begin{itemize}
  \item products and coproducts;
  \item equalizers and coequalizers;
  \item pullbacks and pushouts;
  \item limits and colimits;
  \item the Yoneda embedding;
  \item adjunctions.
\end{itemize}
Each is built from the same template. Recognizing the template is half the
work of understanding the definition.

\subsection{Simplifying: The Uniqueness Schema}

\begin{equation*}
  \forall\, X \in \mathcal{P},\; \exists !\, u : X \to U \;\text{such that}\; \Phi(u).
\end{equation*}

Here $\mathcal{P}$ is the class of competitors, $U$ is the universal object,
and $\Phi(u)$ is the commutativity condition. Reading this as a sentence
in English gives the prose form of any universal property.

\subsection{Formally: Universal Elements via Representable Functors}

\formalnote{Universal properties as natural isomorphisms}
Every universal property can be restated as the assertion that a certain
functor $\mathcal{C} \to \mathbf{Set}$ is \term{representable}: that is,
naturally isomorphic to a hom-functor $\mathrm{Hom}_\mathcal{C}(U, -)$ or
$\mathrm{Hom}_\mathcal{C}(-, U)$. We will not develop $\mathbf{Set}$ or
representability here --- both come in a later chapter --- but it is worth
mentioning that the pattern of this section has a clean abstract
formulation: a universal property is a representation of a certain
functor.


% ═══════════════════════════════════════════════════════════════════════════
\section{Uniqueness Up to Unique Isomorphism}
\label{sec:uniqueness-iso}
% ═══════════════════════════════════════════════════════════════════════════

\subsection{Informally: ``The'' Universal Object}

Universal objects are not, in general, literally unique. There can be many
initial objects in a category. What \emph{is} unique is the
\emph{isomorphism} between any two of them: not just that they are
isomorphic, but that there is exactly one isomorphism connecting them, given
by the universal property itself.

This is the precise sense in which a universal object is ``the'' universal
object. ``The'' here is shorthand for ``unique up to unique isomorphism.''

\subsection{Expanding: Counting Initial Objects}

\begin{example}{Two Initial Objects Must Be Uniquely Isomorphic}
Recall the indiscrete category on $\{A, B\}$. Both $A$ and $B$ are initial.
The unique morphism $A \to B$ (call it $\alpha$) and the unique morphism
$B \to A$ (call it $\beta$) must satisfy:
\begin{align*}
  \beta \circ \alpha &= \mathrm{id}_A
    && \text{(only one morphism $A \to A$, must be the identity)} \\
  \alpha \circ \beta &= \mathrm{id}_B
    && \text{(only one morphism $B \to B$, must be the identity)}.
\end{align*}
So $\alpha$ is an isomorphism with inverse $\beta$. The isomorphism is
unique because $\alpha$ itself was unique (only one morphism $A \to B$
exists).
\end{example}

\subsection{Formalizing: The Theorem}

\textbf{Theorem (Uniqueness of Initial Objects).}\quad If $X$ and $Y$ are
both initial in $\mathcal{C}$, then there is a unique isomorphism
$X \xrightarrow{\sim} Y$.

\textbf{Proof.} Let $f : X \to Y$ and $g : Y \to X$ be the unique morphisms
guaranteed by the initiality of $X$ and $Y$ respectively. Compute:
\begin{align*}
  g \circ f &: X \to X. \\
  \intertext{There is only one morphism $X \to X$ (by initiality of $X$), and
  $\mathrm{id}_X$ is one. So:}
  g \circ f &= \mathrm{id}_X.
  \intertext{Similarly:}
  f \circ g &: Y \to Y, \\
  f \circ g &= \mathrm{id}_Y.
\end{align*}
So $f$ is an isomorphism, with inverse $g$. The isomorphism $f$ is unique
among morphisms $X \to Y$ because $\mathrm{Hom}(X, Y)$ has exactly one
element. \qed

By duality (apply this theorem in $\mathcal{C}^{\mathrm{op}}$):

\textbf{Theorem (Uniqueness of Terminal Objects).}\quad If $X$ and $Y$ are
both terminal in $\mathcal{C}$, then there is a unique isomorphism
$X \xrightarrow{\sim} Y$.

\subsection{Reorganizing: The General Principle}

The argument generalizes. For any universal property, suppose $U$ and $U'$
both satisfy it with the same data. Then:
\begin{itemize}
  \item The universal property of $U'$ applied to $U$ gives a unique
        morphism $u : U \to U'$ compatible with the data.
  \item The universal property of $U$ applied to $U'$ gives a unique
        morphism $u' : U' \to U$.
  \item The composite $u' \circ u : U \to U$ is forced by the universal
        property to be $\mathrm{id}_U$ (uniqueness applied to $U$
        compared against itself).
  \item Similarly $u \circ u' = \mathrm{id}_{U'}$.
\end{itemize}
So $u$ is the unique isomorphism between $U$ and $U'$. This works for every
universal property in this book; the argument is the same in each case.

\subsection{Simplifying: The Slogan}

A universal object is \term{unique up to unique isomorphism}. Or, in
notation:
\begin{equation*}
  U, U' \text{ both satisfy the same universal property}
  \;\;\Longrightarrow\;\;
  \exists !\, u : U \xrightarrow{\sim} U'.
\end{equation*}

\subsection{Formally: Treating Universal Objects as Concrete}

\formalnote{The convention}
We routinely speak of ``the'' initial object, ``the'' terminal object,
``the'' product of $A$ and $B$, and so on, even though strictly speaking
there is a whole equivalence class of such objects. The justification is
exactly the theorem of this section: the choice of representative within
the equivalence class is harmless because any two representatives are
related by a unique canonical isomorphism.

\formalnote{When uniqueness fails}
In some weaker contexts (e.g., bicategories, weak limits), the uniqueness
of the isomorphism relaxes to ``unique up to higher coherent
isomorphism.'' We will encounter such weakened versions in later chapters.
For ordinary (1-)category theory, the strict uniqueness above holds without
modification.


% ═══════════════════════════════════════════════════════════════════════════
\begin{nameorigin}
\term{Product} and \term{coproduct} take their names from set theory:
the Cartesian product $A \times B$ in $\mathbf{Set}$ is the prototype.
The categorical \emph{product} generalizes Cartesian product via the
universal mapping property (maps into $A \times B$ correspond to pairs
of maps into $A$ and $B$). The \emph{coproduct} (also called \emph{sum})
is the dual: in $\mathbf{Set}$ it's the disjoint union; in $\mathbf{Vect}$
it's the direct sum; in $\mathbf{Top}$ it's the disjoint union with
topology. The ``co-'' prefix marks duality (cross-ref Vol I Ch.~2's
opposite categories).
\end{nameorigin}

\section{Toward Products and Coproducts}
\label{sec:toward-products}
% ═══════════════════════════════════════════════════════════════════════════

\subsection{Informally: Two Objects Combined}

Suppose we have two objects $A$ and $B$ in a category, and we want to talk
about a single object that ``contains'' both of them, in some categorical
sense. There are two natural ways to do this:
\begin{itemize}
  \item Combine them so that you can \emph{extract} either one: this leads
        to \emph{products}.
  \item Combine them so that you can \emph{embed} either one: this leads
        to \emph{coproducts}.
\end{itemize}
Each combination has a universal property. This section sketches both as
previews; the full development comes in Chapter~6.

\subsection{Expanding: Pairing and Choice}

\begin{example}{Products as ``Both at Once''}
Imagine you have a Chop step and a Sauté step. A combined task ``Chop and
Sauté'' should let you extract either the chop part or the sauté part.
That is, there should be two ``projections'' $\pi_1 : \mathrm{Chop} \wedge
\mathrm{Saut\acute{e}} \to \mathrm{Chop}$ and $\pi_2 : \mathrm{Chop} \wedge
\mathrm{Saut\acute{e}} \to \mathrm{Saut\acute{e}}$. Given any other step $X$
that can independently be sent to both Chop and Sauté, there should be a
unique way to send $X$ to the combined task.
\end{example}

\begin{example}{Coproducts as ``Either-Or''}
Conversely, a step ``Chop or Sauté'' should accept a chop or a sauté
input and route it through. So there are two ``inclusions'' $\iota_1 :
\mathrm{Chop} \to \mathrm{Chop} \vee \mathrm{Saut\acute{e}}$ and $\iota_2 :
\mathrm{Saut\acute{e}} \to \mathrm{Chop} \vee \mathrm{Saut\acute{e}}$.
Given any other step $Y$ that can receive a chop and can receive a sauté,
there should be a unique way to receive from the disjoint union.
\end{example}

\subsection{Formalizing: The Two Universal Properties}

A \term{product} of $A$ and $B$ in $\mathcal{C}$ is an object $A \times B$
together with morphisms $\pi_1 : A \times B \to A$ and $\pi_2 : A \times
B \to B$, such that:
\begin{quote}
  \emph{for every} object $X$ with morphisms $f : X \to A$ and $g : X \to B$,\\
  \emph{there exists a unique} morphism $\langle f, g \rangle : X \to A
  \times B$ \emph{such that}\\
  $\pi_1 \circ \langle f, g \rangle = f$ \emph{and} $\pi_2 \circ \langle f,
  g \rangle = g$.
\end{quote}
\medskip
\begin{center}
\begin{tikzcd}[column sep=3em, row sep=2.5em]
  & X
    \arrow[dl, "f"', bend right=15]
    \arrow[dr, "g", bend left=15]
    \arrow[d, dashed, "{\langle f, g \rangle}"] & \\
  A & A \times B \arrow[l, "\pi_1"] \arrow[r, "\pi_2"'] & B
\end{tikzcd}
\end{center}
\medskip

A \term{coproduct} of $A$ and $B$ is an object $A + B$ (also written
$A \sqcup B$) with morphisms $\iota_1 : A \to A + B$ and $\iota_2 : B \to
A + B$, such that:
\begin{quote}
  \emph{for every} object $Y$ with morphisms $f : A \to Y$ and $g : B
  \to Y$,\\
  \emph{there exists a unique} morphism $[f, g] : A + B \to Y$
  \emph{such that}\\
  $[f, g] \circ \iota_1 = f$ \emph{and} $[f, g] \circ \iota_2 = g$.
\end{quote}
\medskip
\begin{center}
\begin{tikzcd}[column sep=3em, row sep=2.5em]
  A \arrow[r, "\iota_1"] \arrow[dr, "f"', bend right=15]
  & A + B \arrow[d, dashed, "{[f, g]}"]
  & B \arrow[l, "\iota_2"'] \arrow[dl, "g", bend left=15] \\
  & Y &
\end{tikzcd}
\end{center}

\begin{howtosay}
\begin{itemize}
  \item $A \times B$ — ``$A$ cross $B$,'' ``the product of $A$ and $B$,''
        or sometimes ``$A$ product $B$.'' Some authors write $A \wedge B$
        (read ``$A$ meet $B$'' or ``$A$ and $B$''), especially in
        lattice-theoretic contexts.
  \item $A + B$ — ``$A$ plus $B$,'' ``the coproduct of $A$ and $B$,'' or
        ``$A$ sum $B$.'' The notation $A \sqcup B$ (read ``$A$ disjoint
        union $B$'') is also common; $A \vee B$ (``$A$ join $B$,'' ``$A$ or
        $B$'') appears in lattice-theoretic contexts.
  \item $\pi_1, \pi_2$ — ``pi sub one'' and ``pi sub two,'' or
        ``first/second projection,'' or ``left/right projection.''
  \item $\iota_1, \iota_2$ — ``iota sub one,'' ``iota sub two,'' or
        ``first/second injection'' or ``left/right inclusion.''
  \item $\langle f, g \rangle$ — ``angle $f$ comma $g$,'' ``the pairing of
        $f$ and $g$,'' or ``$f$ paired with $g$.''
  \item $[f, g]$ — ``bracket $f$ comma $g$,'' ``the cotuple of $f$ and
        $g$,'' or ``case analysis on $f$ and $g$.''
\end{itemize}
\end{howtosay}

\subsection{Reorganizing: Duality Again}

A coproduct in $\mathcal{C}$ is a product in $\mathcal{C}^{\mathrm{op}}$.
Every theorem about products has a corresponding theorem about coproducts
obtained by reversing arrows. The pairing $\langle f, g \rangle$ in a
product corresponds to the cotuple $[f, g]$ in a coproduct. The projections
$\pi_i$ correspond to the inclusions $\iota_i$.

We will use this duality without comment for the rest of the book.

\subsection{Simplifying: The Product and Coproduct Schema}

\begin{align*}
  \text{Product:}\quad
    &\forall (X, f, g),\; \exists !\, \langle f, g \rangle : X \to A \times B,\;
    \pi_i \circ \langle f, g \rangle = (f, g)_i. \\
  \text{Coproduct:}\quad
    &\forall (Y, f, g),\; \exists !\, [f, g] : A + B \to Y,\;
    [f, g] \circ \iota_i = (f, g)_i.
\end{align*}

\subsection{Formally: Existence and Notation}

\formalnote{When products exist}
Not every category has products. The categories from Chapter~1 (recipe,
approval chain) generally do not, because there are not enough morphisms to
support pairings. We will exhibit categories that \emph{do} have products
in Chapter~6.

\formalnote{Binary, finite, and infinite}
The product of two objects is the \emph{binary} product. One can also speak
of products of finitely many objects, indexed products
$\prod_{i \in I} A_i$, and infinite products. The universal-property
template extends uniformly to all of these. The same goes for coproducts.

\formalnote{Uniqueness up to unique isomorphism, restated}
By the theorem of \S5.4, any two products $(P, \pi_1, \pi_2)$ and
$(P', \pi'_1, \pi'_2)$ of the same pair $A, B$ are related by a unique
isomorphism $P \xrightarrow{\sim} P'$ commuting with the projections.
Likewise for coproducts. We may therefore speak of ``the'' product and
``the'' coproduct.


% ═══════════════════════════════════════════════════════════════════════════
\section*{Exercises}
\addcontentsline{toc}{section}{Exercises}
% ═══════════════════════════════════════════════════════════════════════════

\begin{exercisesbox}
\begin{qa}
\qaitem{What does it mean for an object $X$ to be initial?}{There is exactly one morphism from $X$ to every object.}
\qaitem{What does it mean for $X$ to be terminal?}{There is exactly one morphism from every object to $X$.}
\qaitem{Is the empty category $\mathbf{0}$ initial or terminal in itself?}{Vacuously its own (trivial) initial object.}
\qaitem{Which object is initial in $\mathbf{2}$?}{The object $0$ (source of the unique non-identity arrow).}
\qaitem{Which is terminal in $\mathbf{2}$?}{The object $1$.}
\qaitem{Can a category have no initial object?}{Yes.}
\qaitem{Can a category have more than one initial object?}{Yes, but any two are uniquely isomorphic.}
\qaitem{Why ``up to unique isomorphism''?}{The isomorphism between two initial objects is itself uniquely determined.}
\qaitem{What is the standard notation for the unique morphism $\mathbf{0} \to A$?}{$!_A$, sometimes called ``the shriek to $A$.''}
\qaitem{What is a zero object?}{An object that is both initial and terminal.}
\qaitem{Given a zero object, what is the zero morphism $A \to B$?}{The composite $A \to \mathbf{0} \to B$ through the zero object.}
\qaitem{What is the signature phrase for a universal property?}{``For every $X$ with such-and-such, there exists a unique $u$ such that some diagram commutes.''}
\qaitem{What four pieces specify a universal property?}{The ambient category, the data, the competitors, and the uniqueness requirement.}
\qaitem{When two objects satisfy the same universal property, how are they related?}{By a unique isomorphism.}
\qaitem{Sketch the uniqueness argument.}{Get morphisms $u : U \to U'$ and $u' : U' \to U$ from each side's universal property; argue $u' \circ u$ and $u \circ u'$ must be identities.}
\qaitem{Why does ``unique morphism $X \to X$'' force $u' \circ u = \mathrm{id}_X$?}{Because $\mathrm{id}_X$ also has the property, and there is only one such morphism.}
\qaitem{What is a product $A \times B$?}{An object equipped with projections $\pi_1, \pi_2$ such that any $(f, g)$ pair factors uniquely through it.}
\qaitem{What is the unique map $X \to A \times B$ induced by $f : X \to A$ and $g : X \to B$?}{$\langle f, g \rangle$, the pairing of $f$ and $g$.}
\qaitem{Which equation does $\langle f, g \rangle$ satisfy?}{$\pi_1 \circ \langle f, g \rangle = f$ and $\pi_2 \circ \langle f, g \rangle = g$.}
\qaitem{What is the coproduct $A + B$?}{An object equipped with inclusions $\iota_1, \iota_2$ such that any $(f, g)$ pair factors uniquely through it.}
\qaitem{What is the unique map $A + B \to Y$ induced by $f : A \to Y$ and $g : B \to Y$?}{$[f, g]$, the cotuple of $f$ and $g$.}
\qaitem{Which equation does $[f, g]$ satisfy?}{$[f, g] \circ \iota_1 = f$ and $[f, g] \circ \iota_2 = g$.}
\qaitem{What is the relation between products and coproducts under duality?}{A coproduct in $\mathcal{C}$ is a product in $\mathcal{C}^{\mathrm{op}}$.}
\qaitem{Why might a category fail to have products?}{Not enough morphisms to support pairings; or no candidate object exists.}
\qaitem{In the recipe category, does $\mathrm{Chop} \times \mathrm{Saut\acute{e}}$ exist?}{Generally not; there's no object with the required pairing.}
\qaitem{If a product exists, is it literally unique?}{No; unique up to unique isomorphism.}
\qaitem{Why is ``unique up to unique iso'' good enough for everyday work?}{We pick a representative and use it; any other representative is interchangeable via the canonical iso.}
\qaitem{What word captures all of products, coproducts, initial, and terminal?}{Universal property (or universal construction).}
\qaitem{How is a universal property reframed in terms of representable functors?}{As a natural isomorphism between some functor and a hom-functor.}
\qaitem{Why bother with universal properties when concrete constructions exist?}{They identify the \emph{essence} of the construction, free of arbitrary internal details.}
\end{qa}
\end{exercisesbox}


% ═══════════════════════════════════════════════════════════════════════════
\section*{Chapter Summary}
\addcontentsline{toc}{section}{Chapter Summary}
% ═══════════════════════════════════════════════════════════════════════════

\begin{summarybox}
An \textbf{initial object} $\mathbf{0}$ has a unique morphism to every
object. A \textbf{terminal object} $\mathbf{1}$ has a unique morphism from
every object. They are dual to each other.

\textbf{Uniqueness up to unique isomorphism}: any two objects satisfying the
same universal property are connected by a unique isomorphism.

The \textbf{universal property pattern}:
\begin{equation*}
  \forall\, X \text{ with data},\; \exists !\, u : X \to U \;\text{or}\; U \to X
  \text{ compatible with the data}.
\end{equation*}

\textbf{Products} ($A \times B$) and \textbf{coproducts} ($A + B$) are the
first universal constructions involving two objects: a product comes with
projections $\pi_1, \pi_2$ and a pairing operation $\langle f, g \rangle$;
a coproduct comes with inclusions $\iota_1, \iota_2$ and a cotuple $[f, g]$.

\medskip
\textbf{Recognizing a universal property:} look for ``for every \ldots,
there exists a unique \ldots such that some diagram commutes.''
\end{summarybox}

% ═══════════════════════════════════════════════════════════════════════════
\section*{Formal Restatement}
\addcontentsline{toc}{section}{Formal Restatement}
% ═══════════════════════════════════════════════════════════════════════════

\begin{formalrestatement}
\textbf{Initial object.}\quad
$\mathbf{0} \in \mathcal{C}$ is initial iff
$|\mathrm{Hom}(\mathbf{0}, A)| = 1$ for every $A \in \mathcal{C}$.

\medskip
\textbf{Terminal object.}\quad
$\mathbf{1} \in \mathcal{C}$ is terminal iff
$|\mathrm{Hom}(A, \mathbf{1})| = 1$ for every $A \in \mathcal{C}$.

\medskip
\textbf{Duality.}\quad
$X$ is terminal in $\mathcal{C}$ iff $X$ is initial in
$\mathcal{C}^{\mathrm{op}}$.

\medskip
\textbf{Uniqueness theorem.}\quad
Any two objects satisfying the same universal property are related by a
unique isomorphism:
\begin{equation*}
  U, U' \text{ both satisfy the property}
  \;\;\Longrightarrow\;\;
  \exists !\, u : U \xrightarrow{\sim} U'.
\end{equation*}

\medskip
\textbf{Product.}\quad
$(A \times B, \pi_1, \pi_2)$ is a product of $A, B$ iff for every
$(X, f, g)$ with $f : X \to A$, $g : X \to B$:
\begin{equation*}
  \exists !\, \langle f, g \rangle : X \to A \times B
  \;\text{with}\;
  \pi_1 \circ \langle f, g \rangle = f,\;
  \pi_2 \circ \langle f, g \rangle = g.
\end{equation*}

\medskip
\textbf{Coproduct.}\quad
$(A + B, \iota_1, \iota_2)$ is a coproduct of $A, B$ iff for every
$(Y, f, g)$ with $f : A \to Y$, $g : B \to Y$:
\begin{equation*}
  \exists !\, [f, g] : A + B \to Y
  \;\text{with}\;
  [f, g] \circ \iota_1 = f,\;
  [f, g] \circ \iota_2 = g.
\end{equation*}

\medskip
\noindent\textit{Chapter~6 develops the full theory of limits and colimits,
of which initial and terminal objects, products and coproducts, are the
first examples.}
\end{formalrestatement}
